% document de test for report-DMKM.cls (LaTeX2e)
% Created May 15, 2012

% We use the style given in rapport-DMKM.cls
\documentclass{report-DMKM}

% this package (optionnal) allows merging rows and columns in latex table 

\usepackage{multirow}

%The following fields (title, author, date, tutors et place) must be present
\title{Place here the title of your internship}
\author{Forename FAMILY\_NAME}
% Defense date
\date{06/29/2012}
\tutors{Supervisors and affiliations}
\place{Location of your internship}
% You can add the logo of the institution where you carry out your internship. Example here ERIC lab

\logo{img/logo-ERIC}{25mm}

\begin{document}

% Cover page (mandatory)
\maketitle

\newpage
% summary in English  (mandatory)
\begin{abstract}
Here is your abstract translated in english.
\end{abstract}

% summary in French  (mandatory)
\begin{resume}
Le pr�sent document est un exemple de fichier \LaTeX utilisant le style \texttt{report-DMKM.cls} qui d�crit la structure d'un rapport de stage pour le Master DMKM, r�dig� en anglais.
Vous devez placer ici votre r�sum� en fran�ais.
Attention, l'en-t�te de votre rapport, comprenant le titre, votre nom, la date, les logos \ldots et les r�sum�s en fran�ais et en anglais, ne doit pas d�passer 1 page au total (il s'agit de la couverture).
\end{resume}

% table of content   (mandatory)
\tabledesmatieres

\section{Hosting institution}

Give an overview about your hosting institution (1 page max)  

\section{Acknowledgement}

You can thank all people and institutions that helped you in your internship

\newpage

\section{Introduction}

This section presents a detailed overview of the work. It gives the objectives/problematics, the methodology you propose and the main results. It ends by a presentation of the content of each section.

The references must be in bibtex format see  
 \texttt{biblio.bib}).
all complementary details are given here  : \url{http://www.bibtex.org/}.
This will give the references like  \cite{HAL02}.

\section{Exemple de section}

You develop the second part of your document .
The final document must be between 20 - 30 pages  , 
If needed, you can add appendices. 
Here bellow a typical plan for this document : 

\begin{enumerate}

\item
Presentation of the subject, expectations etc.

\item
Biblio analysis; this will allow to situate your work within the domain and to take advantage from what it has been done by other authors. 

\item
Your proposal and your justification

\item
results

\item
Discussion.
 
\item
Conclusions and or perspectives
highlight your main results 
\item
references.

\item
Appendices.

\end{enumerate}

 
\newpage

\subsection{Example of subsectionn}

Vous pouvez utiliser des sous-sections.

\subsubsection{example of sub/sub section}

arege  ehheu kieke oe rppf ppjjf je 
you can use more format tools such as  \texttt{paragraph} as followt :

\paragraph{Paragraph 1 : }
blabla blabla blabla

\paragraph{Paragraphe  : }
blabla blabla blabla

\bigskip

To includ  graphics  \texttt{includegraphics} do this :

\begin{center}
\includegraphics[width=2cm]{img/logoECD}
\end{center}
 
Put all your figures in the repository  \texttt{/img},  

\subsection{handling tables}

In \LaTeX , tables are handled by   the command \texttt{tabular} as follow :

\begin{tabular}{|c|c|c|}
% hline draw vertical line
\hline
A & B & C\\
\hline
\hline
$a_1$ & $b_2$ & $c_1$\\
\hline
$a_3$ & $b_2$ & $c_3$\\
\hline
$a_2$ & $b_1$ & $c_3$\\
\hline
\end{tabular}


\medskip

more complex table

\begin{center}
\begin{tabular}{|c||c|c|}
\hline
A & B & C\\
\hline
\hline
\multirow{3}{5mm}{$a_1$}
& $b_2$ & $c_1$\\
% 
\cline{2-3}
& \multicolumn{2}{c|}{$b_2$ et $c_3$}\\
\cline{2-3}
& $b_1$ & $c_3$\\
\hline
\end{tabular}
\end{center}

\subsection{maths}

In \LaTeX, the formula are within  ``dollar'' as follow : $a.x + b = 0$, $x_0$, $y^{12}$, $\sum_{i=0}^{n}{f(i)}$ \ldots
Or :
$$f(x) = \min_{h \in H}{h(\frac{n}{x})}$$
et num�rot�es comme suit :
\begin{center}
\begin{equation}
\label{formule1}
g(z)=\frac{1}{2}f(z) + 1
\end{equation}
\end{center}

this allows equation referencing  \ref{formule1}.

\subsection{to know more about \LaTeX}

 

\begin{itemize}

\item
Joli manual for LaTeX2e de B. Bayart : \url{ftp://ftp.inria.fr/pub/TeX/CTAN/info/JMPL.ps.gz}  
\item
FAQ \LaTeX : \url{http://www.grappa.univ-lille3.fr/FAQ-LaTeX/}

\item
LaTeX2e,   de M. Goosens : \url{http://tex.loria.fr/}

\item
An example of document \LaTeX : \url{http://iml.univ-mrs.fr/~regnier/enseignement/tex-tutorial/}

\item
 tutorial : \url{http://it.metr.ou.edu/latex/}

\end{itemize}

\newpage

\section{Conclusion}

example of citation : \cite{WIT05}

% References are in  biblio.bib   bibtex format
\references

 
\appendixECD

% add first appendix
\annexe{First appendix}
Content
% add second appendix
\annexe{second appendix}
Content

\end{document}
